% vim: ft=tex
% SECTION: SETUP

Building the LED controllers requires some patience, but parts were chosen so
that soldering wouldn't be overly difficult.  This section is a guide for
assembling the board, adding it to a lighting project, and configuring it for
use.

\subsection{Assembling the Board}

First, you need to decide your power configuration.  The LEDs you use will have
some kind of voltage drop; the RGB LEDs I use have roughly a 3V drop per
channel.  So, if I power them with a 5V supply, the LED controller needs to
pick up the 2V slack.  The goal is to minimize the voltage drop across the LED
controller.  The logic runs strictly on 5V only.  You can use two power
supplies (i.e. one for the logic and one for the LEDs) if you'd like; leave
jumper JP1 unconnected.  Or, you can use the same power supply for both the
LEDs; to do that, connect the two pads of JP1 with a piece of wire.

Furthermore, there is a spot for an LM7805 voltage regulator to set the voltage
for the \emph{logic only}.  So, if you want to use a 10V supply source for both
your LEDs and your logic, then you must install an LM7805 in its spot so that
the microcontroller and other devices only get 5V.  However, if you have a 5V
input source, you can tie the JP2 pads together or you can use a piece of wire
to connect pins 1 and 3 of the LM7805 slot.  Do not put an LM7805 in and also
tie the JP2 pads; I'm pretty sure something bad will happen.

The board uses a mixture of surface mount and through-hole components.  I
recommend soldering the surface mount components onto the board first.  My
usual technique for doing so is to tack down one side of a part (e.g. a
resistor) with a bit of solder using tweezers to align the component up with
the other pad.  Then, I go to the other pad and solder it correctly.  Finally,
I go back to the first pad and reflow it to make sure the solder joint is good.
When soldering a component where one of its pads is on the ground plane, I
always tack to the other pad first since the ground plane is hard to heat up.

% TODO: surface mount component description

Next, I recommend tackling the shortest through-hole components first.  The
smaller capacitors would be a good place to start.  Then, solder the SN75176B
(8-pin DIP) chip down.  Next, the four DIP sockets for the microcontroller and
TLC5940s.  Finally, add the electrolytic capacitor, LM7805 (if you are going to
use it), barrel jacks, and RJ-45 jacks.  \textbf{If this is the last board in a
chain} (i.e.  one jack is used for input and the other isn't), solder a
120$\Omega$ resistor to the Rx %TODO
spot on the back side of the board.  This is called a terminating resistor; it
eliminates reflections, which are a major source of communication problems.

Now, you are ready to install the LED drivers and the microcontroller.  First,
make sure that the microcontroller is burned with the bootloader (and that its
fuses are properly set).  You must use an AVR programmer to program the
microcontroller.  Insert the chips into their appropriate sockets.
\textbf{Important:} note that the third TLC5940 (the one that's to the right of
the microcontroller) faces opposite of the other chips -- that is, the first
pin is on the bottom right, not the top left.

\subsection{Setting up the Power Supply}

I use standard 2.54mm barrel jacks for the power input of the board.  Depending
on what power supply you use, you may need to convert its power jacks to
support this scheme.

\subsection{Installing the LEDs}

Next, you'll need to wire your LEDs to the board.  PLEDVCC holes provide
positive voltage from the LED power supply.  PLED holes connect to the current
sink outputs of the LED drivers.  You can solder wires from the LEDs directly
to the board or you can probably find some kind of spade plugs small enough to
fit in the holes if you want to be able to disconnect the LEDs from the board.
Figure~\ref{fig:setup:leds} shows a typical setup.

\begin{rfigure*}
    \label{fig:intro:nhpls}
    \includegraphics[scale=1.0]{fig/photo.pdf}
    %\includegraphics[scale=0.6]{fig/fig-intro-nhpls.pdf}
    % TODO nhpls board photo
    \caption{LED controller board used in the Next House Party Lighting System}
\end{rfigure*}


